% =============================================================================
% Chapter: Decision Trees & Classification
% =============================================================================

\chapter{Decision Trees \& Classification}

\section{Εισαγωγή}

Τα \keyword{Decision Trees} είναι supervised learning models που χρησιμοποιούνται για classification και regression. Χτίζουν ένα δέντρο αποφάσεων βασισμένο στα features των δεδομένων.

\subsection{Βασική Δομή}

\begin{center}
\begin{tikzpicture}[
    level distance=1.8cm,
    level 1/.style={sibling distance=4cm},
    level 2/.style={sibling distance=2cm},
    edge from parent/.style={draw, arrow}
]
    \node[treenode, fill=ntua-blue!30, font=\bfseries] {Root}
        child {
            node[treenode] {Internal}
            child {node[leafnode] {Class A}}
            child {node[leafnode] {Class B}}
        }
        child {
            node[treenode] {Internal}
            child {node[leafnode] {Class A}}
        };
\end{tikzpicture}
\end{center}

\begin{definitionbox}[Βασικοί Όροι]
\begin{itemize}
    \item \textbf{Root Node}: Αρχικός κόμβος (feature με μεγαλύτερη διαχωριστική ικανότητα)
    \item \textbf{Internal Node}: Ενδιάμεσος κόμβος (test σε κάποιο attribute)
    \item \textbf{Leaf Node}: Τελικός κόμβος (class label/prediction)
    \item \textbf{Branch}: Σύνδεση μεταξύ κόμβων (outcome του test)
\end{itemize}
\end{definitionbox}

% =============================================================================
\section{Βασικές Έννοιες \& Τύποι}

\subsection{Entropy (Εντροπία)}

Μετράει την \keyword{αβεβαιότητα} ή \keyword{impurity} ενός συνόλου δεδομένων.

\begin{formulabox}[Entropy]
\begin{equation}
H(S) = -\sum_{i=1}^{c} p_i \cdot \log_2(p_i)
\end{equation}

όπου $p_i$ = πιθανότητα της κλάσης $i$ στο σύνολο $S$, και $c$ = αριθμός κλάσεων
\end{formulabox}

\begin{center}
\begin{tikzpicture}
    \begin{axis}[
        width=10cm,
        height=6cm,
        xlabel={$p$ (πιθανότητα class 1)},
        ylabel={Entropy $H$},
        domain=0.01:0.99,
        samples=100,
        grid=both,
        grid style={line width=.1pt, draw=gray!30},
        axis lines=middle,
        ymin=0, ymax=1.1,
        xmin=0, xmax=1,
        title={\textbf{Entropy για Binary Classification}},
        title style={font=\small}
    ]
        \addplot[thick, secondary, smooth] {-x*ln(x)/ln(2) - (1-x)*ln(1-x)/ln(2)};
        \node[pin={[pin edge={secondary, thick}]above right:{Max at $p=0.5$}}, secondary]
            at (axis cs:0.5,1) {};
    \end{axis}
\end{tikzpicture}
\end{center}

\begin{infobox}[Ιδιότητες Entropy]
\begin{itemize}
    \item $H(S) = 0$ → \textbf{Τέλεια καθαρότητα} (όλα της ίδιας κλάσης)
    \item $H(S) = 1$ → \textbf{Μέγιστη αβεβαιότητα} (για 2 κλάσεις: 50-50 split)
    \item $H(S) \in [0, \log_2(c)]$ όπου $c$ = αριθμός κλάσεων
\end{itemize}
\end{infobox}

\begin{examplebox}[Υπολογισμός Entropy]
Dataset $S$: 9 ``Yes'', 5 ``No'' (Total: 14 samples)

\begin{align*}
p(\text{Yes}) &= \frac{9}{14} = 0.643 \\
p(\text{No}) &= \frac{5}{14} = 0.357 \\[10pt]
H(S) &= -\left(\frac{9}{14}\right)\log_2\left(\frac{9}{14}\right) - \left(\frac{5}{14}\right)\log_2\left(\frac{5}{14}\right) \\
&= -(0.643)(-0.637) - (0.357)(-1.485) \\
&= 0.410 + 0.530 = \boxed{0.940 \text{ bits}}
\end{align*}
\end{examplebox}

% =============================================================================
\subsection{Information Gain}

Μετράει τη \keyword{μείωση της αβεβαιότητας} όταν διαχωρίζουμε με βάση ένα attribute.

\begin{formulabox}[Information Gain]
\begin{equation}
IG(S, A) = H(S) - \sum_{v \in \text{Values}(A)} \frac{|S_v|}{|S|} \cdot H(S_v)
\end{equation}

\begin{itemize}
    \item $H(S)$ = entropy πριν το split
    \item $S_v$ = subset όπου attribute $A = v$
    \item $|S_v|/|S|$ = weighted proportion
\end{itemize}
\end{formulabox}

\begin{warningbox}[Κανόνας Επιλογής]
Διαλέγουμε το attribute με το \textbf{ΜΕΓΑΛΥΤΕΡΟ Information Gain}!
\end{warningbox}

% =============================================================================
\subsection{Gain Ratio (C4.5)}

Το Information Gain \important{ευνοεί attributes με πολλές τιμές}. Το Gain Ratio το διορθώνει.

\begin{formulabox}[Gain Ratio]
\begin{equation}
\text{GainRatio}(S, A) = \frac{IG(S, A)}{\text{SplitInfo}(S, A)}
\end{equation}

\begin{equation}
\text{SplitInfo}(S, A) = -\sum_{v \in \text{Values}(A)} \frac{|S_v|}{|S|} \cdot \log_2\left(\frac{|S_v|}{|S|}\right)
\end{equation}
\end{formulabox}

% =============================================================================
\subsection{Gini Index}

Εναλλακτικό μέτρο impurity (χρησιμοποιείται στο CART).

\begin{formulabox}[Gini Index]
\begin{equation}
\text{Gini}(S) = 1 - \sum_{i=1}^{c} p_i^2
\end{equation}

\begin{itemize}
    \item $\text{Gini} = 0$ → pure (μία κλάση)
    \item $\text{Gini} = 0.5$ → max impurity (2 κλάσεις 50-50)
\end{itemize}
\end{formulabox}

% =============================================================================
\section{Αλγόριθμοι}

\subsection{ID3 Algorithm}

\begin{algorithmbox}[ID3 (Iterative Dichotomiser 3)]
\begin{lstlisting}[language=pseudo]
function ID3(S, Attributes):
    if all samples in S belong to same class C:
        return Leaf(C)

    if Attributes is empty:
        return Leaf(majority class in S)

    A_best = argmax IG(S, A) for A in Attributes

    Create Node for A_best

    for each value v of A_best:
        S_v = samples where A_best = v
        if S_v is empty:
            Add Leaf(majority class of S)
        else:
            Add ID3(S_v, Attributes - {A_best})

    return Node
\end{lstlisting}
\end{algorithmbox}

\subsection{Σύγκριση Αλγορίθμων}

\begin{center}
\begin{tikzpicture}
    \matrix (m) [
        matrix of nodes,
        nodes={
            draw=primary,
            minimum width=2.8cm,
            minimum height=1cm,
            anchor=center,
            font=\small
        },
        column 1/.style={nodes={fill=ntua-blue!20, font=\bfseries}},
        row 1/.style={nodes={fill=secondary!30, font=\bfseries}},
        column sep=-\pgflinewidth,
        row sep=-\pgflinewidth
    ] {
        Algorithm & Criterion & Splits & Special \\
        ID3 & Information Gain & Multi-way & Categorical only \\
        C4.5 & Gain Ratio & Multi-way & Handles missing \\
        CART & Gini Index & Binary only & Regression support \\
    };
\end{tikzpicture}
\end{center}

% =============================================================================
\section{Worked Example: Κατασκευή Δέντρου}

\begin{examplebox}[Play Tennis Dataset]
Κατασκευάζουμε decision tree για το Play Tennis dataset.
\end{examplebox}

\begin{center}
\begin{tabular}{cccc|c}
\toprule
\textbf{Outlook} & \textbf{Temp} & \textbf{Humidity} & \textbf{Wind} & \textbf{Play?} \\
\midrule
Sunny & Hot & High & Weak & No \\
Sunny & Hot & High & Strong & No \\
Overcast & Hot & High & Weak & Yes \\
Rain & Mild & High & Weak & Yes \\
Rain & Cool & Normal & Weak & Yes \\
Rain & Cool & Normal & Strong & No \\
\bottomrule
\end{tabular}
\end{center}

\textbf{Βήμα 1: Υπολογισμός $H(S)$}
\begin{align*}
H(S) &= -\frac{9}{14}\log_2\frac{9}{14} - \frac{5}{14}\log_2\frac{5}{14} = 0.940 \text{ bits}
\end{align*}

\textbf{Βήμα 2: Υπολογισμός IG για κάθε attribute}

\begin{center}
\begin{tikzpicture}
    \begin{axis}[
        ybar,
        width=10cm,
        height=5cm,
        bar width=1.2cm,
        ylabel={Information Gain},
        symbolic x coords={Outlook, Humidity, Wind, Temperature},
        xtick=data,
        nodes near coords,
        nodes near coords align={vertical},
        ymin=0,
        ymax=0.35,
        title={\textbf{Information Gain για κάθε Attribute}}
    ]
        \addplot[fill=secondary!70] coordinates {
            (Outlook, 0.246)
            (Humidity, 0.151)
            (Wind, 0.048)
            (Temperature, 0.029)
        };
    \end{axis}
\end{tikzpicture}
\end{center}

\textbf{Βήμα 3: Επιλογή Root} → \keyword{Outlook} (μέγιστο IG = 0.246)

\textbf{Τελικό Decision Tree:}

\begin{center}
\begin{tikzpicture}[
    level distance=2cm,
    level 1/.style={sibling distance=4cm},
    level 2/.style={sibling distance=2cm},
    edge from parent/.style={draw, ->, >=Stealth, thick}
]
    \node[treenode, fill=ntua-blue!30, minimum size=1.2cm, font=\bfseries] {Outlook}
        child {
            node[treenode, minimum size=1cm] {Humidity}
            child {node[leafnode, fill=accent!30] {No} edge from parent node[left, font=\scriptsize] {High}}
            child {node[leafnode, fill=success!30] {Yes} edge from parent node[right, font=\scriptsize] {Normal}}
            edge from parent node[left, font=\small] {Sunny}
        }
        child {
            node[leafnode, fill=success!30, minimum height=0.8cm] {Yes}
            edge from parent node[left, font=\small] {Overcast}
        }
        child {
            node[treenode, minimum size=1cm] {Wind}
            child {node[leafnode, fill=success!30] {Yes} edge from parent node[left, font=\scriptsize] {Weak}}
            child {node[leafnode, fill=accent!30] {No} edge from parent node[right, font=\scriptsize] {Strong}}
            edge from parent node[right, font=\small] {Rain}
        };
\end{tikzpicture}
\end{center}

% =============================================================================
\section{Pruning (Κλάδεμα)}

\begin{center}
\begin{tikzpicture}
    % Pre-pruning
    \node[draw=success, fill=success!20, rounded corners=10pt,
          minimum width=5cm, minimum height=2.5cm] (pre) at (0,0) {};
    \node[anchor=north, font=\bfseries\color{success}] at (pre.north) {Pre-Pruning};
    \node[anchor=center, align=center, font=\small] at (pre.center) {
        Σταματάει νωρίς\\[3pt]
        • Max depth\\
        • Min samples per leaf\\
        • Min IG threshold
    };

    % Post-pruning
    \node[draw=secondary, fill=secondary!20, rounded corners=10pt,
          minimum width=5cm, minimum height=2.5cm] (post) at (7,0) {};
    \node[anchor=north, font=\bfseries\color{secondary}] at (post.north) {Post-Pruning};
    \node[anchor=center, align=center, font=\small] at (post.center) {
        Χτίζει πλήρες,\\
        μετά κόβει\\[3pt]
        • Reduced Error\\
        • Cost Complexity
    };

    % Arrow
    \draw[->, >=Stealth, thick, dashed] (pre.east) -- (post.west)
        node[midway, above, font=\small] {vs};
\end{tikzpicture}
\end{center}

% =============================================================================
\section{Συχνά Λάθη}

\begin{warningbox}[Κοινά Λάθη στις Εξετάσεις]
\begin{enumerate}
    \item \textbf{Λάθος Entropy}: $H = p \cdot \log_2(p)$ \\
          \textcolor{success}{Σωστό}: $H = \mathbf{-}p \cdot \log_2(p)$

    \item \textbf{Weighted Average}: $H(S|A) = \frac{H(S_1) + H(S_2)}{2}$ \\
          \textcolor{success}{Σωστό}: $H(S|A) = \frac{|S_1|}{|S|}H(S_1) + \frac{|S_2|}{|S|}H(S_2)$

    \item \textbf{Σύγχυση ID3/C4.5}:
          \begin{itemize}
              \item ID3 → Information Gain
              \item C4.5 → Gain Ratio
          \end{itemize}
\end{enumerate}
\end{warningbox}

% =============================================================================
\section{Σύνοψη - Key Points}

\begin{center}
\begin{tikzpicture}
    % Summary boxes
    \node[draw=ntua-blue, fill=ntua-blue!10, rounded corners=5pt,
          minimum width=14cm, minimum height=8cm] (main) at (0,0) {};

    % Formulas section
    \node[anchor=north west, font=\bfseries\Large\color{ntua-blue}]
        at ([xshift=5pt, yshift=-5pt]main.north west) {Τύποι (ΑΠΟΜΝΗΜΟΝΕΥΣΕ!)};

    \node[anchor=north west, align=left] at ([xshift=10pt, yshift=-30pt]main.north west) {
        \begin{tabular}{ll}
            \textbf{Entropy} & $H(S) = -\sum p_i \log_2(p_i)$ \\[5pt]
            \textbf{Info Gain} & $IG = H(S) - H(S|A)$ \\[5pt]
            \textbf{Gain Ratio} & $GR = IG / \text{SplitInfo}$ \\[5pt]
            \textbf{Gini} & $G(S) = 1 - \sum p_i^2$ \\
        \end{tabular}
    };

    % Algorithms section
    \node[anchor=west, align=left] at ([xshift=10pt, yshift=-50pt]main.west) {
        \begin{tabular}{lll}
            \textbf{ID3} & Information Gain & Multi-way splits \\
            \textbf{C4.5} & Gain Ratio & Handles continuous \\
            \textbf{CART} & Gini Index & Binary splits only \\
        \end{tabular}
    };

    % Key properties
    \node[anchor=south west, align=left, font=\small]
        at ([xshift=10pt, yshift=10pt]main.south west) {
        \textcolor{success}{$H = 0$} → Pure node |
        \textcolor{accent}{$H = 1$} → Max uncertainty |
        \textcolor{secondary}{Bigger IG = Better split}
    };
\end{tikzpicture}
\end{center}
